\section{INTRODUÇÃO}
\label{sec:intro}

A qualidade da educação básica é um dos principais determinantes do desenvolvimento econômico de longo prazo, já que as habilidades cognitivas da população afetam o crescimento \parencite{hanushek2012}. No Brasil, o Índice de Desenvolvimento da Educação Básica (IDEB) monitora o ensino público combinando fluxo escolar e proficiência em língua portuguesa e matemática numa medida comparável entre municípios, estados e redes de ensino \parencite{fernandes2007}. Apesar dos avanços das últimas duas décadas, estados do Norte e do Nordeste seguem abaixo da média nacional, reflexo de desigualdades estruturais \parencite{inep2025}.

O Maranhão está entre os casos mais críticos. Com a menor renda domiciliar \textit{per capita} do país no Censo Demográfico 2022 e indicadores educacionais historicamente abaixo da média nacional, o estado concentra municípios em alta vulnerabilidade socioeconômica \parencite{ibge2022, inep2025}. Ainda assim, o desempenho varia entre municípios de condições semelhantes, o que coloca uma pergunta direta: o que diferencia, dentro de um mesmo patamar de vulnerabilidade, os municípios de maior IDEB daqueles de menor IDEB?

A literatura mostra que o nível socioeconômico é o fator mais associado ao desempenho, mas que as escolas também fazem diferença \parencite{coleman1966, brooke2008}, por meio da qualificação docente \parencite{rivkin2005} e da infraestrutura \parencite{alves2013, vasconcelos2021}. Análises do Maranhão que isolem fatores intracontextuais, isto é, aqueles que se mantêm associados ao desempenho mesmo entre municípios de contexto socioeconômico semelhante, são escassas. Algoritmos de agrupamento \parencite{romero2020, kalita2025} permitem segmentar esses contextos antes das comparações.

O objetivo geral deste artigo é identificar, entre municípios maranhenses de contexto socioeconômico semelhante, os fatores escolares associados ao desempenho superior no IDEB, considerando os 216 dos 217 municípios com IDEB 2025 divulgado (anos iniciais do ensino fundamental, rede municipal). Os objetivos específicos são: agrupar os municípios por contexto socioeconômico, com o algoritmo \textit{K-Means} aplicado à renda \textit{per capita} e à proporção de moradores em pobreza extrema do Censo Demográfico 2022; comparar, dentro de cada grupo, os municípios de alto e de baixo desempenho quanto à infraestrutura, à qualificação docente e ao fluxo escolar registrados no Censo Escolar 2025, com o teste de Mann-Whitney, tamanho de efeito e correção para comparações múltiplas; e discutir as implicações dos achados para as políticas educacionais municipais.

Diferentemente dos modelos globais, que comparam municípios de contextos muito distintos, o estudo compara municípios dentro de grupos socioeconômicos semelhantes, oferece evidência sobre um estado pouco estudado e propõe uma abordagem replicável em outras unidades federativas. Seguem-se o referencial teórico, os procedimentos metodológicos (fontes, variáveis e os dois estágios da análise), os resultados e sua discussão e as considerações finais, com a síntese dos achados e as limitações.
